\section{Conclusão e Trabalhos Futuros}

% SCRIPT: Resumir a conclusão central. O projeto provou ser uma alternativa viável.
% Destacar que as decisões de software garantiram a estabilidade sob alta carga.
% Frisamos a delimitação: é ponte de aquisição, não para controle crítico.
\begin{frame}{Síntese do Trabalho}
    \begin{itemize}
        \item \textbf{Objetivo Cumprido:} O WCAN é uma ponte de aquisição viável em bancada, eliminando chicotes longos e mantendo a transparência lógica do barramento CAN.
        \item \textbf{A Recomendação Final:} A arquitetura mais promissora para o veículo é o \textbf{\textit{Multicast} aliado ao agrupamento temporal (\textit{batching})}.
        \item \textbf{Delimitação Operacional:} A solução preserva a organização do protocolo, mas \textbf{não substitui} enlaces físicos para as funções de controle crítico do protótipo.
    \end{itemize}
\end{frame}

% SCRIPT: Apontar as próximas etapas.
% Responder à pergunta clássica: "Testou no carro?". Não, porque a base teórica e de protocolo (saturação de filas, etc) precisava ser garantida antes de introduzir as variáveis físicas (EMI, vibração).
\begin{frame}{Limitações e Próximos Passos}
    \begin{alertblock}{O Sistema foi testado no carro?}
        \textbf{Não.} O escopo do TCC foi provar a solidez do protocolo, das filas no RTOS e do transporte em ESP-NOW. Levar ao veículo antes de mapear a saturação de \textit{software} (em 1000 Hz) mascararia as falhas com problemas de vibração ou interferência eletromagnética (EMI).
    \end{alertblock}

    \vspace{0.4cm}
    \textbf{Trabalhos Futuros:}
    \begin{itemize}
        \item \textbf{Validação Física:} Instalação no veículo para ensaios de autonomia, blindagem e fixação.
        \item \textbf{Integração Automotiva:} Uso da rede \textit{TWAI} para reinjeção e consumo contínuo de mensagens no barramento elétrico do carro.
        \item \textbf{Comunicação Bidirecional:} Expandir testes para nós que produzem e consomem dados simultaneamente no WCAN.
    \end{itemize}
\end{frame}
